%************************************************
\chapter{Methods} \label{ch:methods}
%************************************************

% TODO chapter opener (Mizzaro: objective + ordered section preview) — draft once
% the sections exist.

\section{Aim of the Work} \label{sec:methods:aim}

\acp{SAE} have become the standard tool for decomposing the activations of
\iac{LLM} into a dictionary of concept-like features
(\cref{sec:background:ae:sae}). Their
reach is bounded, though, by a structural trade-off between sparsity and
reconstruction \citep{gaoScalingEvaluatingSparse2024}: because each concept
enters the reconstruction through a single scalar coefficient, reconstructing an
activation faithfully tends to require many concepts active at once, whereas the
interpretability that motivates the decomposition rests on keeping only a few
active per token \citep{hubenSparseAutoencodersFind2023,%
brickenMonosemanticityDecomposingLanguage2023}.
Pushed to a useful sparsity, the reconstruction degrades far enough that the
reconstructed activation leaves the data manifold. This is more than an
aesthetic flaw: a downstream circuit then processes an out-of-distribution
signal even before any intervention, and when a feature is amplified or ablated
the resulting change in behaviour conflates the intended edit with this
uncontrolled reconstruction error, so the very analyses the decomposition is
built for become less reliable than they appear.

The limitation lies in the representation rather than the training: a scalar can
record only how strongly a concept is present. This thesis asks whether giving
each concept a \textit{vector} embedding relaxes the trade-off. The idea is
familiar from \acp{CEM}
(\cref{sec:background:xai:c-xai}), which
replace the scalar concept bottleneck of \iac{CBM} with per-concept embeddings
and so recover the accuracy a narrow bottleneck gives up
\citep{espinosazarlengaConceptEmbeddingModels2022} --- but \acp{CEM} are
supervised, trained against concept labels. The aim here is to carry the
vector-concept idea into the \textit{unsupervised} setting of activation
decomposition, and to do so from first principles, defining a probabilistic
generative model whose inference procedure and training objective follow from
the model itself rather than being assembled by hand. The starting point is the
\ac{LCBM} of \citet{desantisBetterGeneralizationInterpretability2025}, an
unsupervised concept-based image classifier whose encoder produces a vector
embedding per concept, gates each concept by comparing that embedding to a
learnable prototype, and reconstructs the input and simultaneously classifies it
from the embeddings of the active concepts. \ac{LCBM} supplied the idea rather
than the architecture: this thesis keeps its core structure --- concepts as
embeddings gated against prototypes and decoded by reconstruction --- but
rebuilds it from the ground up as a fully variational generative model, giving
the concept embeddings a probabilistic latent that \ac{LCBM} leaves
deterministic and dropping the supervised classification task, so that the model
reconstructs activations from no labels at all.

That model --- the \ac{VAEE} --- sits where three of the families surveyed in
\cref{ch:background} meet. From
\acp{CEM} it takes the representation of a concept as a high-dimensional
embedding rather than a scalar; from \acp{VAE} it takes a probabilistic latent
and amortised variational inference
(\cref{sec:background:ae:vae});
and from \acp{VQVAE} it takes a learned set of concept prototypes together with
a discrete choice of which concepts are active\graffito{sure about this?}
(\cref{sec:background:ae:vae:arch}). The
result is a probabilistic counterpart to discrete quantisation: independent
Bernoulli activations route continuous, prototype-anchored embeddings into
separate concept slots, retaining the expressivity of a continuous latent while
keeping the routing sparse and selective. Nothing in this construction is
specific to language models --- the vector-concept primitive applies wherever a
scalar concept code trades capacity for simplicity --- but this thesis develops
and demonstrates it for one application: the decomposition of \ac{LLM}
activations.

What this thesis contributes, then, is the architecture itself and its
derivation: an unsupervised, vector-valued concept model obtained from a
probabilistic generative model rather than assembled by hand. Its aim is
correspondingly bounded --- to expand the representational capacity of a concept
code, leaving the interpretability and steering of the recovered concepts to
future work --- and \cref{ch:experiments} tests it on
the decomposition of \ac{LLM} activations.

\section{Definitions/Notation} \label{sec:methods:def}
% TODO: the reconciled VAEE notation (see .claude/notation.md "VAEE notation
% LOCKED") + the "prototype" disambiguation (VAEE h_k vs XAI input-instance).

\section{Model Development} \label{sec:methods:model}
    \subsection{Probabilistic Formulation} \label{sec:methods:model:proba}
    % TODO: Gated Gaussian generative model (c -> z -> x) + PGM; amortised
    % inference (encoder, embedding posterior, prototype-similarity gating,
    % Gumbel-Sigmoid, gating step); the ELBO (3 terms) -> design choices ->
    % practical loss; similarity metrics. Full A.1/A.2 algebra goes to
    % Chapters/C0A_VAEEDerivations (app:vaee:gating, app:vaee:elbo).

    \subsection{Architectures} \label{sec:methods:model:arch}
    % TODO: encoder/decoder instantiations (shallow/linear/MLP); stacked vs summed
    % topologies (stacked = headline); interventions (formulation supports them;
    % evaluation deferred to future work); decoder orthogonality (formulated then
    % dropped as empirically negligible).

\section*{Conclusions}
% TODO: bridge to C04 (Experiments).


%*****************************************
%*****************************************
%*****************************************
%*****************************************
%*****************************************
